\chapter{Projeto do Sistema}
\label{capitulo:projeto}

Este capítulo apresenta os principais elementos de projeto do simulador \textit{Royal Legacy}. São detalhadas a arquitetura do sistema, os protótipos das interfaces, a modelagem de dados para persistência, os diagramas estruturais e comportamentais baseados na Unified Modeling Language (UML), bem como a representação lógica dos algoritmos centrais de validação e comunicação.

As decisões de projeto aqui expostas visam atender aos requisitos de modularidade, customização visual e execução assíncrona da Inteligência Artificial, conforme justificado na seção \ref{secao:justificativa}.

\section{Arquitetura do projeto}

A arquitetura do \textit{Royal Legacy} foi projetada seguindo o princípio da separação de responsabilidades, dividindo o sistema em camadas distintas que se comunicam através de protocolos padronizados. Essa abordagem é fundamental para garantir que o processamento pesado da IA não impacte a fluidez da interface gráfica \cite{Fowler2003}.

A Figura \ref{figura:arquitetura_sistema} apresenta o diagrama de implantação que ilustra como as partes do sistema foram divididas e como elas se comunicam.

%\figura{caminho}{tamanho cm}{legenda}{referencia}
\figura{./capitulos/projeto/arquitetura_implantação.png}{12cm}{Diagrama de Implantação da arquitetura do Royal Legacy.}{figura:arquitetura_sistema}

O fluxo de comunicação segue o modelo Cliente-Servidor local:
\begin{enumerate}
	\item \textbf{Camada de Apresentação (Godot Engine):} Atua como o cliente (Frontend). Gerencia a interface, captura entradas do usuário e envia o estado do jogo (formato FEN) para o intermediário.
	\item \textbf{Camada de Comunicação (\textit{Middleware} Python):} Um \textit{script} que roda em segundo plano, agindo como ponte. Ele orquestra a execução da IA e traduz os comandos.
	\item \textbf{Camada Lógica (Stockfish Binary):} O motor de xadrez (Backend). Recebe configurações via protocolo UCI e devolve os cálculos matemáticos.
\end{enumerate}

\section{Protótipos}
\label{secao:prototipos}

Nesta seção são apresentados os protótipos de baixa fidelidade das principais interfaces do \textit{Royal Legacy}. Estes rascunhos serviram como base para o desenvolvimento visual, focando na disposição dos elementos e na Experiência do Usuário (UX), não representando as telas finais com texturas e cores aplicadas.

\subsection{Menu Principal e Opções}

O protótipo do Menu Principal (Figura \ref{figura:prototipo_menu}) foi desenhado para ser minimalista, oferecendo acesso rápido aos modos de jogo e configurações. O protótipo da tela de Opções (Figura \ref{figura:prototipo_opcoes}) detalha os rascunhos dos seletores para volume, nível de dificuldade da IA (Fácil, Médio, Difícil) e a galeria de temas visuais.

\figura{./capitulos/projeto/prototipo_menu.png}{10cm}{Rascunho da interface do Menu Principal.}{figura:prototipo_menu}
\figura{./capitulos/projeto/prototipo_opcoes.png}{10cm}{Rascunho da interface de configurações de áudio, dificuldade e temas.}{figura:prototipo_opcoes}

\subsection{Tabuleiro de Jogo}

O protótipo da tela de jogo (Figura \ref{figura:prototipo_jogo}) apresenta a grade $8\times8$ centralizada, ladeada por rascunhos dos painéis de peças capturadas, cronômetro (opcional) e o botão para retornar ao menu.

\figura{./capitulos/projeto/prototipo_jogo.png}{10cm}{Rascunho da interface principal do tabuleiro e painéis laterais.}{figura:prototipo_jogo}

\section{Diagrama Entidade Relacionamento (DER)}

Dado que o \textit{Royal Legacy} é um simulador local focado na jogabilidade e não na gestão complexa de usuários ou grandes volumes de dados relacionais, optou-se por não utilizar um Banco de Dados Relacional tradicional (como PostgreSQL ou MySQL). Em vez disso, a persistência de dados é realizada através de arquivos de configuração locais (formato JSON ou dicionários serializados na Godot), o que simplifica a arquitetura e reduz o consumo de recursos da máquina \cite{Date2003}.

No entanto, para representar como as estruturas de dados necessárias para o funcionamento do sistema foram modeladas de forma lógica e normalizada, apresenta-se na Figura \ref{figura:der_sistema} o Diagrama Entidade Relacionamento (DER) que rege a persistência das configurações e estados de partidas.

\figura{./capitulos/projeto/der_jogo.png}{12cm}{Diagrama Entidade Relacionamento para persistência local.}{figura:der_sistema}

O DER segue as normas formais de modelagem:
\begin{itemize}
	\item \textbf{Configuracao\_Local:} Armazena as preferências do usuário (Volume, Dificuldade, ID\_Tema\_Ativo).
	\item \textbf{Tema\_Visual:} Estrutura os caminhos dos arquivos (\textit{sprites}) e códigos de cores para a customização dinâmica da interface.
	\item \textbf{Historico\_Partida (Opcional):} Permite salvar o estado atual do tabuleiro em formato FEN para retomada posterior.
\end{itemize}

Todos os campos visíveis nos protótipos da seção \ref{secao:prototipos} (como seletores de volume e temas) possuem representação correspondente neste modelo de dados.

\section{Diagrama de Casos de Uso}

O Diagrama de Casos de Uso (Figura \ref{figura:caso_uso}) descreve as funcionalidades do sistema sob a perspectiva dos atores envolvidos. No contexto do \textit{Royal Legacy}, o ator principal é o **Jogador Humano**, enquanto o **Motor Stockfish** atua como um ator secundário invocado pelo sistema nos modos PvE.

\figura{./capitulos/projeto/diagrama_uso.png}{12cm}{Diagrama de Casos de Uso do simulador.}{figura:caso_uso}

Os principais casos de uso incluem:
\begin{itemize}
	\item \textbf{Configurar Jogo:** Permite alterar volume, dificuldade e tema visual.
		\item \textbf{Iniciar Partida:** Subdividido em Jogador vs Jogador (PvP) ou Jogador vs Máquina (PvE).
			\item \textbf{Realizar Jogada:** Inclui a validação de regras lógicas.
				\item \textbf{Solicitar Melhor Movimento:** Caso de uso executado pelo sistema que invoca o ator Stockfish no modo PvE.
				\end{itemize}
				
				\section{Diagrama de Classes}
				
				A organização estrutural do código na Godot Engine, embora baseada em Cenas e Nós, pode ser representada formalmente através de um Diagrama de Classes UML. A Figura \ref{figura:diagrama_classes} apresenta as principais classes de domínio do \textit{Royal Legacy} e suas relações.
				
				\figura{./capitulos/projeto/diagrama_classes.png}{14cm}{Diagrama de Classes detalhando a estrutura GDScript e middleware Python.}{figura:diagrama_classes}
				
				Destacam-se as seguintes classes e relações:
				\begin{itemize}
					\item \textbf{StructuralManager:} Atua como o nó raiz (\texttt{estruturador.tscn}), gerenciando os estados de troca de cenas (Menus e Tabuleiro).
					\item \textbf{ChessBoard:} Classe central que detém a matriz do tabuleiro, o histórico de lances e a lógica de validação de regras.
					\item \textbf{PythonBridge:** Classe responsável por executar o processo externo via \textit{script} Python de forma assíncrona.
						\item \textbf{AudioManager (Singleton):** Classe Autoload (\texttt{tela\_sons.gd}) para execução contínua de áudio.
						\end{itemize}
						
						\section{Algoritmos propostos}
						
						Em conformidade com as diretrizes deste trabalho, esta seção não apresenta códigos-fonte brutos (GDScript ou Python), mas sim a representação lógica dos algoritmos que constituem contribuições claras para a solução do problema de pesquisa.
						
						O Algoritmo \ref{alg:solicitacao_assincrona} descreve a lógica implementada na Godot para solicitar a jogada da Inteligência Artificial sem travar a interface visual, utilizando a execução de subprocessos externos.
						
						\begin{algorithm}
							\caption{Solicitação Assíncrona de Jogada da IA (Godot Side).}
							\label{alg:solicitacao_assincrona}
							
							\Inicio{
								$estado\_fen$ \gets $obter\_estado\_atual\_tabuleiro()$\;
								$dificuldade$ \gets $obter\_configuracao\_dificuldade()$\;
								
								// Preparar o Middleware Python como processo externo\\
								$comando$ \gets $"python chess\_bridge.py --fen " + estado\_fen + " --diff " + dificuldade$\;
								
								// Executar de forma assíncrona (thread separada)\\
								\Se{ $OS.execute\_assincrono(comando, callback\_funcao\_retorno)$ }
								{
									$bloquear\_entrada\_usuario\_humano()$\;
									$exibir\_icone\_processando()$\;
								}
							}
						\end{algorithm}
						
						O Algoritmo \ref{alg:middleware_python} representa a lógica central do \textit{middleware} Python, que gerencia a comunicação via \textit{pipes} de terminal com o binário do Stockfish utilizando o protocolo UCI.
						
						\begin{algorithm}
							\caption{Middleware Python de Comunicação com Stockfish (UCI Protocol).}
							\label{alg:middleware_python}
							
							\Entrada{$arg\_fen$ (String), $arg\_diff$ (String)}
							\Saida{Melhor jogada (String formato "e2e4")}
							
							\Inicio{
								// Abrir processo binário do Stockfish com pipes de leitura/escrita\\
								$stockfish$ \gets $subprocess.Popen("./stockfish", stdin=PIPE, stdout=PIPE)$ \cite{Stockfish2023}\;
								
								// Configurar dificuldade e posição via UCI\\
								$enviar\_comando\_uci(stockfish, "setoption name Skill Level value " + arg\_diff)$ \;
								$enviar\_comando\_uci(stockfish, "position fen " + arg\_fen)$ \;
								
								// Solicitar cálculo restringindo tempo ou profundidade\\
								$enviar\_comando\_uci(stockfish, "go movetime 2000")$\;
								
								// Ler saída do Stockfish aguardando a linha "bestmove"\\
								$linha$ \gets $ler\_saida\_pipe(stockfish)$ \;
								\Enquanto{ $"bestmove" \notin linha $}
								{
									$linha$ \gets $ler\_saida\_pipe(stockfish)$ \;
								}
								
								// Extrair e retornar apenas as coordenadas da jogada\\
								$jogada\_final$ \gets $extrair\_coordenadas(linha)$ \;
								$fechar\_processo(stockfish)$ \;
								\Retorna{$jogada\_final$}
							}
						\end{algorithm}